\documentclass[11pt,a4paper]{article}
\usepackage[T1]{fontenc}
\usepackage[utf8]{inputenc}
\usepackage[brazilian]{babel}
\usepackage{lmodern}
\usepackage[margin=2.3cm]{geometry}
\usepackage{enumitem}
\usepackage{xcolor}
\usepackage{fancyhdr}

\pagestyle{fancy}
\fancyhf{}
\rhead{\small Pesquisa acadêmica, UNISINOS}
\cfoot{\small\thepage}
\renewcommand{\headrulewidth}{0pt}

\setlength{\parindent}{0pt}
\setlength{\parskip}{0.55em}
\setlist[enumerate]{leftmargin=*, itemsep=0.25em, topsep=0.2em}
\setlist[itemize]{leftmargin=*, itemsep=0.2em, topsep=0.2em}
\newcommand{\tema}[1]{\subsection*{\normalsize #1}}

\begin{document}

\begin{center}
  {\Large\bfseries Cardápios digitais em bares e restaurantes de Porto Alegre}\\[0.3em]
  {\large Convite para entrevista com gestores}
\end{center}

\section*{\normalsize Sobre a pesquisa}

Sou César Hoffmann, aluno de Ciência da Computação na UNISINOS, e estou concluindo meu Trabalho de Conclusão de Curso, orientado pelo Prof.\ Dr.\ Guilherme Lacerda. A pesquisa investiga como bares e restaurantes de Porto Alegre lidam com cardápios, impressos ou digitais, e com o atendimento desde 2020, e como os clientes vivenciam essas escolhas.

Até agora, ouvimos mais de cem frequentadores de bares e restaurantes e gestores de casas com perfis diferentes. Cada estabelecimento fez escolhas próprias, e todas elas interessam à pesquisa: não existe resposta certa ou errada. Gostaríamos muito de ouvir a experiência de vocês.

\section*{\normalsize Como seria a conversa}

\begin{itemize}
  \item Cerca de 30 minutos, presencialmente ou por videoconferência, no dia e horário que for melhor para vocês.
  \item Gravada apenas para transcrição, com o seu consentimento. A gravação não é divulgada.
  \item O nome da casa só aparece no trabalho se vocês autorizarem; caso contrário, usamos uma descrição genérica. Seu nome pessoal não é divulgado.
  \item Não é preciso preparar nada nem compartilhar números sigilosos. Qualquer pergunta pode ser pulada.
  \item Ao final, envio um resumo dos resultados da pesquisa.
\end{itemize}

\section*{\normalsize Temas e perguntas}

As perguntas abaixo servem de guia. A conversa é livre e pode seguir outros caminhos.

\tema{1. A casa}
\begin{enumerate}
  \item Como e quando a casa surgiu? Que tipo de cliente vocês atendem?
  \item Qual é o seu papel no estabelecimento?
  \item Como funciona hoje o atendimento: como o cliente consulta o cardápio, faz o pedido e paga?
\end{enumerate}

\tema{2. O cardápio e o atendimento}
\begin{enumerate}[resume]
  \item Como vocês decidiram o formato de cardápio e de atendimento que usam hoje?
  \item Esse formato mudou ao longo do tempo, por exemplo durante ou depois da pandemia? Por quê?
  \item O que pesou nessa decisão (custos, equipe, clientes, legislação, outra coisa)?
\end{enumerate}

\tema{3. Sistemas e operação}
\begin{enumerate}[resume]
  \item Que sistemas a casa usa no dia a dia (caixa, pedidos, cozinha, delivery)? Foram contratados ou desenvolvidos?
  \item Como o pedido chega à cozinha? E como vocês atualizam preços e itens que acabaram?
  \item O que acontece quando a internet ou algum sistema para de funcionar?
\end{enumerate}

\tema{4. Clientes}
\begin{enumerate}[resume]
  \item Que retorno vocês recebem dos clientes sobre o cardápio e o atendimento?
  \item Algum grupo de clientes tem mais dificuldade, por exemplo pessoas idosas ou com baixa visão? Como vocês lidam com isso?
\end{enumerate}

\tema{5. Futuro}
\begin{enumerate}[resume]
  \item Vocês pensam em mudar algo no cardápio ou no atendimento nos próximos anos?
  \item Se pudessem pedir uma coisa a quem desenvolve sistemas para restaurantes, o que seria?
\end{enumerate}

\vspace{0.6em}
\textbf{Contato:} César Hoffmann \textbullet{} cesarhoff@unisinos.br \textbullet{} (51) 99804-4896 (WhatsApp)\\
\textbf{Orientador:} Prof.\ Dr.\ Guilherme Lacerda, Escola Politécnica, UNISINOS

\end{document}
